% Results Section: Semantic Structure of LLM Routing

\subsection{Semantic Structure of Task Difficulty}
\label{sec:semantic_structure}

Before presenting our routing algorithm, we first establish that LLM routing is a \emph{structured} 
problem amenable to learning. Figure~\ref{fig:lmsys_holdout_structure} visualizes 1,871 prompts from 
the LMSYS dev and holdout sets projected onto the first two principal components of a 32-component PCA model 
trained on sentence embeddings (all-MiniLM-L6-v2). Despite the 2D projection capturing only 5.39\% 
of the total variance (PC1: 3.10\%, PC2: 2.29\%), clear semantic structure emerges that provides 
the empirical foundation for cost-effective routing.

\paragraph{Bimodal Semantic Distribution.} As illustrated in Figure~\ref{fig:lmsys_holdout_structure}, 
the semantic task structure of the LMSYS dev and holdout sets (N=1,871) is inherently bimodal. The spatial 
distribution reveals two distinct clusters separated by a decision boundary at PC1 $\approx$ 0.3:

\begin{itemize}[noitemsep,topsep=2pt]
    \item \textbf{Low PC1 Cluster} (82.4\%, 1,541 prompts): Represents semantic regions where the 
    performance delta between flagship and mid-tier models is marginal---routine conversational queries, 
    simple creative tasks, and straightforward information requests.
    \item \textbf{High PC1 Cluster} (17.6\%, 330 prompts): Isolates complex coding challenges, 
    multi-step reasoning tasks, and advanced mathematical proofs that require flagship model capabilities.
\end{itemize}

\noindent This spatial separation provides the empirical justification for contextual routing: a static 
strategy would either over-spend on the 82.4\% of routine tasks or under-perform on the 17.6\% of 
critical tasks. By exploiting the PC1 decision boundary, our routing approach targets a theoretical 
80\%+ reduction in operational costs without compromising quality on the high-variance reasoning cluster.

\paragraph{Statistical Significance of Low Variance.} While 5.39\% captured variance may appear modest, 
in a high-dimensional embedding space ($d=384$), the top two components isolate the primary ``task complexity'' 
axis. Even at low variance ratios, these clusters are statistically distinct and provide sufficient signal 
for the 1.26$\times$ near-optimal recovery observed in our Hybrid model (Section~\ref{sec:results}). 
The kernel density estimation (KDE) contours in Figure~\ref{fig:lmsys_holdout_structure} confirm that 
the Low PC1 cluster exhibits strong spatial coherence, indicating these are dominant features of the 
traffic distribution rather than noise.

\paragraph{Implications for Routing.} This structured distribution has three key implications:

\textbf{(1) Proof of Learnability.} The presence of bimodal semantic clusters demonstrates that task 
difficulty is not random but predictable from prompt embeddings. A routing policy can generalize from 
training examples to unseen prompts by exploiting this structure. The clear separation visible even 
in 2D projection provides immediate evidence that routing is a solvable learning problem.

\textbf{(2) Economic Opportunity.} The 82.4\% / 17.6\% split reveals that the majority of production 
traffic occupies semantic regions where expensive flagship models are unnecessary. This asymmetry 
creates a natural target for cost optimization: by routing the Low PC1 cluster to mid-tier models 
while preserving flagship quality for the High PC1 cluster, we can achieve substantial operational 
savings without sacrificing user experience on critical tasks.

\textbf{(3) Cold-start Priors are Viable.} The bimodal structure enables Bayesian recalibration: 
if warmup priors assign initial beliefs about task difficulty based on semantic neighborhoods (e.g., 
PC1 position), these priors can be efficiently refined through bandit feedback. This motivates our 
hybrid approach (Section~\ref{sec:method}) that combines structural priors with adaptive learning.

\paragraph{Validation with Unseen Data.} Critically, Figure~\ref{fig:lmsys_holdout_structure} uses 
the LMSYS \emph{dev and holdout} sets---prompts representing production-realistic evaluation data. 
These sets provide a cleaner, more realistic signal of the LMSYS task distribution compared to 
aggregated training distributions used for model development. The clear bimodal structure in this 
evaluation data demonstrates that semantic task organization is robust and generalizable, providing 
confidence that routing policies can transfer to production deployment where similar distributional 
characteristics are expected.

\paragraph{Defense Against Common Critiques.} Two potential concerns merit explicit address:

\textbf{``Is 5.39\% variance enough?''} In high-dimensional spaces, low absolute variance can still 
represent dominant structure. Our results demonstrate that even this modest 2D projection achieves 
strong cluster separation and enables near-optimal routing performance. The full 32-component PCA 
(35.14\% variance) provides even richer signals for the routing algorithm, but the 2D visualization 
suffices to prove the existence of learnable structure.

\textbf{``Why use dev/holdout sets instead of the full 80K training data?''} The dev and holdout 
sets represent production-realistic evaluation data, avoiding the mixed signals often found in 
large aggregated training distributions. By demonstrating bimodal structure on evaluation-quality 
prompts, we provide stronger evidence of the distributional characteristics relevant for production 
deployment where the router must handle diverse queries in real-time. The N=1,871 sample provides 
sufficient statistical power to establish bimodality while maintaining data quality.

